% GENERATED FILE. Edit canonical lessons, then run npm run generate:books.
% canonical-source-hash: fnv1a64:6eded5956b7610ec
% canonical-lessons: JA-C346-kondo, JA-C346-konkai, JA-C346-zenkai, JA-C346-jikai, JA-C346-jiki, JA-R346-review-341-346

\chapter{This Time and Next Time}
\label{ch:ja-this-time-and-next-time}
\hlchaptermodality{346}

\begin{chapteropening}
\textbf{By the end of this chapter:} \emph{I can talk about this time, last time and next time.}

Complete the last lesson of chapter 346: I can talk about this time, last time and next time.
\end{chapteropening}

\section[\texorpdfstring{\ja{こんど} (kondo)}{こんど (kondo)}]{\ja{こんど} (kondo) — next time; this time}

\label{lesson:JA-C346-kondo}



\begin{quote}
Before the new one: say the Japanese for to shift a time, then the Japanese for to bring forward.
\end{quote}

\subsection*{You'll want to know: \ja{こんど}}
\textbf{\ja{こんど}} — \emph{kondo} — \textquotedblleft{}next time; this time\textquotedblright{}.

\begin{quote}
\textbf{\ja{こんど} \ja{じかんが} \ja{あったら、うちに} \ja{きて} \ja{ください。}} — \emph{kondo jikan ga attara, uchi ni kite kudasai} — \textquotedblleft{}Next time you have a moment, please come to my house.\textquotedblright{}
\end{quote}

\begin{cousinweb}
\textbf{\ja{こんど}} (\emph{kondo}) is Sino-Japanese (kango): two Chinese characters read with their on'yomi, readings borrowed from Chinese. They are \emph{kon}, \textquotedblleft{}now\textquotedblright{} (Middle Chinese \emph{kim}) and \emph{do}, \textquotedblleft{}degrees\textquotedblright{} (\emph{duH}) (KANJIDIC2; Middle Chinese in Baxter's transcription).
\end{cousinweb}

\begin{grammarlens}[title={What you've built}]
One more word for saying what happens if.
\end{grammarlens}

\subsection*{Your turn}
\begin{itemize}
\raggedright
  \item \emph{Say it:} \emph{kondo}
  \item \emph{Say it:} \emph{kondo}, once more
  \item \emph{Say it:} \emph{kondo jikan ga attara, uchi ni kite kudasai} — Next time you have a moment, please come to my house.
  \item \emph{Recall:} say the Japanese for to shift a time, then the Japanese for to bring forward, then say \emph{kondo} again
\end{itemize}

\begin{tcolorbox}[breakable,colback=teal!4,colframe=teal!35!black,title={Before you move on}]
What does \ja{こんど} mean? (\textquotedblleft{}Next time; this time\textquotedblright{}.) Say it once more.
\end{tcolorbox}
\section[\texorpdfstring{\ja{こんかい} (konkai)}{こんかい (konkai)}]{\ja{こんかい} (konkai) — this time; on this occasion}

\label{lesson:JA-C346-konkai}



\begin{quote}
Before the new one: say the Japanese for to bring forward, then the Japanese for next time.
\end{quote}

\subsection*{You'll want to know: \ja{こんかい}}
\textbf{\ja{こんかい}} — \emph{konkai} — \textquotedblleft{}this time; on this occasion\textquotedblright{}.

\begin{quote}
\textbf{\ja{こんかい} \ja{まけたら、つぎは} \ja{がんばります。}} — \emph{konkai maketara, tsugi wa ganbarimasu} — \textquotedblleft{}If we lose this time, we will try harder next time.\textquotedblright{}
\end{quote}

\begin{cousinweb}
\textbf{\ja{こんかい}} (\emph{konkai}) is Sino-Japanese (kango): two Chinese characters read with their on'yomi, readings borrowed from Chinese. They are \emph{kon}, \textquotedblleft{}now\textquotedblright{} (Middle Chinese \emph{kim}) and \emph{kai}, \textquotedblleft{}-times\textquotedblright{} (\emph{hwoj}) (KANJIDIC2; Middle Chinese in Baxter's transcription).
\end{cousinweb}

\begin{grammarlens}[title={What you've built}]
One more word for saying what happens if.
\end{grammarlens}

\subsection*{Your turn}
\begin{itemize}
\raggedright
  \item \emph{Say it:} \emph{konkai}
  \item \emph{Say it:} \emph{konkai}, once more
  \item \emph{Say it:} \emph{konkai maketara, tsugi wa ganbarimasu} — If we lose this time, we will try harder next time.
  \item \emph{Recall:} say the Japanese for to bring forward, then the Japanese for next time, then say \emph{konkai} again
\end{itemize}

\begin{tcolorbox}[breakable,colback=teal!4,colframe=teal!35!black,title={Before you move on}]
What does \ja{こんかい} mean? (\textquotedblleft{}This time; on this occasion\textquotedblright{}.) Say it once more.
\end{tcolorbox}
\section[\texorpdfstring{\ja{ぜんかい} (zenkai)}{ぜんかい (zenkai)}]{\ja{ぜんかい} (zenkai) — last time; the previous time}

\label{lesson:JA-C346-zenkai}



\begin{quote}
Before the new one: say the Japanese for next time, then the Japanese for on this occasion.
\end{quote}

\subsection*{You'll want to know: \ja{ぜんかい}}
\textbf{\ja{ぜんかい}} — \emph{zenkai} — \textquotedblleft{}last time; the previous time\textquotedblright{}.

\begin{quote}
\textbf{\ja{ぜんかいより} \ja{はやく} \ja{ついたら、えきで} \ja{まって} \ja{ください。}} — \emph{zenkai yori hayaku tsuitara, eki de matte kudasai} — \textquotedblleft{}If you arrive earlier than last time, please wait at the station.\textquotedblright{}
\end{quote}

\begin{cousinweb}
\textbf{\ja{ぜんかい}} (\emph{zenkai}) is Sino-Japanese (kango): two Chinese characters read with their on'yomi, readings borrowed from Chinese. They are \emph{zen}, \textquotedblleft{}before\textquotedblright{} (Middle Chinese \emph{dzen}) and \emph{kai}, \textquotedblleft{}-times\textquotedblright{} (\emph{hwoj}) (KANJIDIC2; Middle Chinese in Baxter's transcription).
\end{cousinweb}

\begin{grammarlens}[title={What you've built}]
One more word for saying what happens if.
\end{grammarlens}

\subsection*{Your turn}
\begin{itemize}
\raggedright
  \item \emph{Say it:} \emph{zenkai}
  \item \emph{Say it:} \emph{zenkai}, once more
  \item \emph{Say it:} \emph{zenkai yori hayaku tsuitara, eki de matte kudasai} — If you arrive earlier than last time, please wait at the station.
  \item \emph{Recall:} say the Japanese for next time, then the Japanese for on this occasion, then say \emph{zenkai} again
\end{itemize}

\begin{tcolorbox}[breakable,colback=teal!4,colframe=teal!35!black,title={Before you move on}]
What does \ja{ぜんかい} mean? (\textquotedblleft{}Last time; the previous time\textquotedblright{}.) Say it once more.
\end{tcolorbox}
\section[\texorpdfstring{\ja{じかい} (jikai)}{じかい (jikai)}]{\ja{じかい} (jikai) — the next occasion}

\label{lesson:JA-C346-jikai}



\begin{quote}
Before the new one: say the Japanese for on this occasion, then the Japanese for the previous time.
\end{quote}

\subsection*{You'll want to know: \ja{じかい}}
\textbf{\ja{じかい}} — \emph{jikai} — \textquotedblleft{}the next occasion\textquotedblright{}.

\begin{quote}
\textbf{\ja{じかいの} \ja{かいぎは、てんきに} \ja{よります。}} — \emph{jikai no kaigi wa, tenki ni yorimasu} — \textquotedblleft{}Whether the next meeting goes ahead depends on the weather.\textquotedblright{}
\end{quote}

\begin{cousinweb}
\textbf{\ja{じかい}} (\emph{jikai}) is Sino-Japanese (kango): two Chinese characters read with their on'yomi, readings borrowed from Chinese. They are \emph{ji}, \textquotedblleft{}next\textquotedblright{} (Middle Chinese \emph{tshijH}) and \emph{kai}, \textquotedblleft{}-times\textquotedblright{} (\emph{hwoj}) (KANJIDIC2; Middle Chinese in Baxter's transcription).
\end{cousinweb}

\begin{grammarlens}[title={What you've built}]
One more word for saying what happens if.
\end{grammarlens}

\subsection*{Your turn}
\begin{itemize}
\raggedright
  \item \emph{Say it:} \emph{jikai}
  \item \emph{Say it:} \emph{jikai}, once more
  \item \emph{Say it:} \emph{jikai no kaigi wa, tenki ni yorimasu} — Whether the next meeting goes ahead depends on the weather.
  \item \emph{Recall:} say the Japanese for on this occasion, then the Japanese for the previous time, then say \emph{jikai} again
\end{itemize}

\begin{tcolorbox}[breakable,colback=teal!4,colframe=teal!35!black,title={Before you move on}]
What does \ja{じかい} mean? (\textquotedblleft{}The next occasion\textquotedblright{}.) Say it once more.
\end{tcolorbox}
\section[\texorpdfstring{\ja{じき} (jiki)}{じき (jiki)}]{\ja{じき} (jiki) — a time; a season; a period}

\label{lesson:JA-C346-jiki}



\begin{quote}
Before the new one: say the Japanese for the previous time, then the Japanese for the next occasion.
\end{quote}

\subsection*{You'll want to know: \ja{じき}}
\textbf{\ja{じき}} — \emph{jiki} — \textquotedblleft{}a time; a season; a period\textquotedblright{}.

\begin{quote}
\textbf{\ja{あつい} \ja{じきに} \ja{なったら、みずを} \ja{のんで} \ja{ください。}} — \emph{atsui jiki ni nattara, mizu o nonde kudasai} — \textquotedblleft{}When the hot season comes, please drink water.\textquotedblright{}
\end{quote}

\begin{cousinweb}
\textbf{\ja{じき}} (\emph{jiki}) is Sino-Japanese (kango): two Chinese characters read with their on'yomi, readings borrowed from Chinese. They are \emph{ji}, \textquotedblleft{}time\textquotedblright{} (Middle Chinese \emph{dzyi}) and \emph{ki}, \textquotedblleft{}period\textquotedblright{} (\emph{gi}) (KANJIDIC2; Middle Chinese in Baxter's transcription).
\end{cousinweb}

\begin{grammarlens}[title={What you've built}]
One more word for saying what happens if.
\end{grammarlens}

\subsection*{Your turn}
\begin{itemize}
\raggedright
  \item \emph{Say it:} \emph{jiki}
  \item \emph{Say it:} \emph{jiki}, once more
  \item \emph{Say it:} \emph{atsui jiki ni nattara, mizu o nonde kudasai} — When the hot season comes, please drink water.
  \item \emph{Recall:} say the Japanese for the previous time, then the Japanese for the next occasion, then say \emph{jiki} again
\end{itemize}

\begin{tcolorbox}[breakable,colback=teal!4,colframe=teal!35!black,title={Before you move on}]
What does \ja{じき} mean? (\textquotedblleft{}A time; a season; a period\textquotedblright{}.) Say it once more.
\end{tcolorbox}
\section[(dialogue)]{Looking back: from Persuading Someone to This Time and Next Time}

\label{lesson:JA-R346-review-341-346}



\begin{quote}
Say the Japanese for to persuade and to urge.
\end{quote}

\subsection*{Your turn}
Say these aloud:

\begin{itemize}
\raggedright
  \item to persuade, to urge, to ask someone to stay, as much as possible, for example
  \item a cause, an influence, a result, the reason is, an attitude
  \item to trust, confidence, to stand up for, to be found out, to jump to conclusions
  \item a detailed plan, to get things ready, an errand, a schedule, going smoothly
  \item to change a plan, to postpone, to cancel, to shift a time, to bring forward
  \item next time, on this occasion, the previous time, the next occasion, a period of time
\end{itemize}

\begin{tcolorbox}[breakable,colback=teal!4,colframe=teal!35!black,title={Before you move on}]
To persuade? (\textbf{\ja{せっとくする}}, \emph{settoku suru}.) To urge? (\textbf{\ja{うながす}}, \emph{unagasu}.) To ask someone to stay? (\textbf{\ja{ひきとめる}}, \emph{hikitomeru}.) As much as possible? (\textbf{\ja{なるべく}}, \emph{narubeku}.) For example? (\textbf{\ja{たとえば}}, \emph{tatoeba}.) A cause? (\textbf{\ja{げんいん}}, \emph{gen'in}.) An influence? (\textbf{\ja{えいきょう}}, \emph{eikyō}.) A result? (\textbf{\ja{けっか}}, \emph{kekka}.) The reason is? (\textbf{\ja{なぜなら}}, \emph{nazenara}.) An attitude? (\textbf{\ja{たいど}}, \emph{taido}.) To trust? (\textbf{\ja{しんようする}}, \emph{shin'yō suru}.) Confidence? (\textbf{\ja{しんらい}}, \emph{shinrai}.) To stand up for? (\textbf{\ja{かばう}}, \emph{kabau}.) To be found out? (\textbf{\ja{ばれる}}, \emph{bareru}.) To jump to conclusions? (\textbf{\ja{きめつける}}, \emph{kimetsukeru}.) A detailed plan? (\textbf{\ja{けいかく}}, \emph{keikaku}.) To get things ready? (\textbf{\ja{よういする}}, \emph{yōi suru}.) An errand? (\textbf{\ja{ようじ}}, \emph{yōji}.) A schedule? (\textbf{\ja{にってい}}, \emph{nittei}.) Going smoothly? (\textbf{\ja{じゅんちょう}}, \emph{junchō}.) To change a plan? (\textbf{\ja{へんこうする}}, \emph{henkō suru}.) To postpone? (\textbf{\ja{えんきする}}, \emph{enki suru}.) To cancel? (\textbf{\ja{とりけす}}, \emph{torikesu}.) To shift a time? (\textbf{\ja{ずらす}}, \emph{zurasu}.) To bring forward? (\textbf{\ja{はやめる}}, \emph{hayameru}.) Next time? (\textbf{\ja{こんど}}, \emph{kondo}.) On this occasion? (\textbf{\ja{こんかい}}, \emph{konkai}.) The previous time? (\textbf{\ja{ぜんかい}}, \emph{zenkai}.) The next occasion? (\textbf{\ja{じかい}}, \emph{jikai}.) A period of time? (\textbf{\ja{じき}}, \emph{jiki}.)
\end{tcolorbox}
